% !TEX root = ../main.tex

% 专硕请取消下面注释
% \makeatletter
%   \def\hitsz@csubjecttitle{类别}
% \makeatother

\hitszsetup{
  %******************************
  % 注意：
  %   1. 配置里面不要出现空行
  %   2. 不需要的配置信息可以删除
  %******************************
  %
  %=====
  % 秘级
  %=====
  statesecrets={公开},
  natclassifiedindex={TM391.1},
  intclassifiedindex={004.8},
  %
  %=========
  % 中文信息
  %=========
  % ctitleone={基于结构与语义分析的检索增强大语言模型生成方法研究},%本科生封面使用
  % ctitletwo={关键技术的研究},%本科生封面使用
  ctitlecover={基于结构与语义分析的大语言模型检索增强生成方法},%放在封面中使用，自由断行
  ctitle={基于结构与语义分析的大语言模型检索增强生成方法},%放在原创性声明中使用
  % csubtitle={一条副标题}, %一般情况没有，可以注释掉
  cxueke={电子信息},
  cpostgraduatetype={专业},
  csubject={电子信息},
  % csubject={机械工程},
  caffil={计算机科学与技术学院},
  caffil={哈尔滨工业大学（深圳）},
  cauthor={李士朋},
  csupervisor={汤步洲 教授},
  % cassosupervisor={某某某 教授}, % 副导师
  % ccosupervisor={某某某 教授}, % 合作导师
  % 日期自动使用当前时间，若需指定按如下方式修改：
  cdate={2026年5月},
  % 指定第二页封面的日期，即答辩日期
  cdatesecond={2026年5月},
  cstudentid={23S151134},
  % cstudenttype={同等学力人员}, %非全日制教育申请学位者
  %（同等学力人员）、（工程硕士）、（工商管理硕士）、
  %（高级管理人员工商管理硕士）、（公共管理硕士）、（中职教师）、（高校教师）等
  %
  %
  %=========
  % 英文信息
  %=========
  etitle={Large Language Model Retrieval-Augmented Generation Methods Based on Structural and Semantic Analysis},
  % esubtitle={This is the sub title},
  exueke={Electronic Information},
  esubject={Electronic Information},
  eaffil={Harbin Institute of Technology, Shenzhen},
  eauthor={Shipeng Li},
  esupervisor={Prof. Buzhou Tang},
  % eassosupervisor={XXX},
  % ecosupervisor={Prof. XXX}, % Co-Supervisor off Campus
  % 日期自动生成，若需指定按如下方式修改：
  edate={May, 2026},
  % estudenttype={Master of Engineering},
  %
  % 关键词用“英文逗号”分割
  ckeywords={检索增强生成, 大语言模型, 结构化推理, 联合语义建模},
  ekeywords={Retrieval-Augmented Generation, Large Language Models, Structured Reasoning, Joint Semantic Modeling},
}

% 中文摘要
\begin{cabstract}

  随着大语言模型（Large Language Models, LLMs）的快速发展，其在自然语言处理等多个领域中表现出良好的性能。
  然而，由于蕴含在模型参数中知识的静态性，LLMs 容易因知识过时而产生“幻觉”现象。
  检索增强生成（Retrieval-Augmented Generation, RAG）通过引入外部知识增强了 LLMs 生成结果的事实性与可靠性，
  但现有的 RAG 方法在知识密集型多跳推理任务等复杂场景下仍存在上下文噪声较大、语义表达能力不足以及推理能力有限等问题。
  针对上述问题，本文围绕“结构分析”与“语义分析”两个方向展开研究，提出了两种改进方法。

  针对复杂多跳问答任务中上下文信噪比较低的问题，提出了一种基于结构分析的检索增强生成方法。
  通过从检索文档中抽取短语结构三元组，构建结构化推理路径，并基于推理路径精确索引相关文本，
  构建“结构引导+文本补充”的增强上下文表示，以提高知识密度，降低上下文噪声。
  在三个典型多跳问答数据集 MuSiQue、2WikiMultiHopQA 和 HotpotQA 上的实验结果表明，
  所提出的方法在各项指标上均优于对比方法。
  其中，在 2WikiMultiHopQA 数据集上，EM 和 F1 分别达到 55.20\% 和 65.09\%，
  相较次优方法分别提升 2.80\% 和 3.76\%。
  在所有数据集上的平均性能达到 55.20\%，取得最优整体表现。
  同时，在保证性能的前提下，所提出的方法将平均输入 Token 长度减少约 70\%--80\%，
  显著提升了推理效率与结果可解释性。

  针对基于图结构的 RAG 方法依赖复杂关系建模且构建成本较高的问题，
  提出了一种基于结构与语义联合建模的检索增强生成方法。
  在短语结构三元组和文本片段的基础上，
  引入实体语义信息抽取，
  并提出基于句子加权的段落表示方法，
  构建“结构 + 文本 + 语义”联合增强上下文，
  提升上下文的信息密度与表达能力。
  在 MuSiQue、2WikiMultiHopQA 和 HotpotQA 数据集上进行实验，
  所提出的方法的包含匹配准确率（Contain-Acc.）和
  大语言模型评估准确率（GPT-Acc.）分别达到 58.57\% 和 57.93\%，
  在整体性能上优于多种基准方法。
  所提出的方法在保持较低检索成本的同时，有效提升了大语言模型的推理能力与生成性能。

\end{cabstract}

% 英文摘要
\begin{eabstract}

  With the rapid development of Large Language Models (LLMs), they have demonstrated strong performance in many fields such as natural language processing.
  However, due to the static nature of knowledge stored in model parameters, LLMs are prone to generating hallucinations caused by outdated knowledge.
  Retrieval-Augmented Generation (RAG) enhances the factuality and reliability of LLM-generated results by introducing external knowledge.
  Nevertheless, existing RAG methods still suffer from issues such as excessive contextual noise, insufficient semantic representation capability, and limited reasoning ability in complex scenarios including knowledge-intensive multi-hop reasoning tasks.
  To address these challenges, this thesis conducts research from two perspectives: structural analysis and semantic analysis, and proposes two improved methods.

  To address the low signal-to-noise ratio of context in complex multi-hop question answering tasks, this thesis proposes a retrieval-augmented generation method based on structural analysis.
  By extracting phrase-structured triples from retrieved documents, the method constructs structured reasoning paths and accurately indexes relevant texts based on these reasoning paths.
  A “structure-guided + text-supplemented” enhanced context representation is then constructed to improve knowledge density and reduce contextual noise.
  Experimental results on three representative multi-hop question answering benchmarks, namely MuSiQue, 2WikiMultiHopQA, and HotpotQA, demonstrate that the proposed method consistently outperforms competing approaches across all evaluation metrics.
  Specifically, on the 2WikiMultiHopQA dataset, the proposed method achieves an EM score of 55.20\% and an F1 score of 65.09\%, surpassing the second-best method by 2.80\% and 3.76\%, respectively.
  The proposed method achieves an average overall performance of 55.20\% across all datasets, obtaining the best comprehensive performance.
  Meanwhile, while maintaining strong performance, the proposed method reduces the average input token length by approximately 70\%--80\%, significantly improving reasoning efficiency and result interpretability.

  To address the high construction cost and complex relational modeling dependency of graph-based RAG methods, this thesis further proposes a retrieval-augmented generation method based on joint structural and semantic modeling.
  Built upon phrase-structured triples and text passages, the method further introduces entity-level semantic information extraction and proposes a sentence-weighted paragraph representation approach.
  A jointly enhanced context integrating “structure + text + semantics” is then constructed to improve context knowledge density and semantic representation capability.
  Experimental results on MuSiQue, 2WikiMultiHopQA, and HotpotQA show that the proposed method achieves a Containment Matching Accuracy (Contain-Acc.) of 58.57\% and a Large Language Model-based Evaluation Accuracy (GPT-Acc.) of 57.93\%, outperforming multiple baseline methods in overall performance.
  While maintaining relatively low retrieval cost, the proposed method effectively improves the reasoning ability and generation performance of large language models.
  
\end{eabstract}
